% Generated draft: requires mathematical and source review.
\documentclass[11pt]{article}
\usepackage[margin=25mm]{geometry}
\usepackage{fontspec}
\setmainfont{Latin Modern Roman}
\usepackage{amsmath,amssymb,mathtools,unicode-math}
\setmathfont{Latin Modern Math}
\usepackage{enumitem,needspace,xurl,hyperref}
\setlength{\parindent}{0pt}
\setlength{\parskip}{5pt}
\setlength{\emergencystretch}{3em}
\newcommand{\N}{\mathbb{N}}
\newcommand{\Z}{\mathbb{Z}}
\newcommand{\Q}{\mathbb{Q}}
\newcommand{\R}{\mathbb{R}}
\newcommand{\C}{\mathbb{C}}
\newcommand{\F}{\mathbb{F}}
\newcommand{\A}{\mathbb{A}}
\newcommand{\PP}{\mathbb{P}}
\newcommand{\E}{\mathbb{E}}
\newcommand{\eps}{\varepsilon}
\newcommand{\dd}{\,\mathrm{d}}
\begin{document}
\section*{UnsolvedMath ID: 2306034}
\textbf{Source problem code:} catalogue entry
\subsection{Definitions and mathematical statement}
Let $\mathcal A$ consist of analytic functions $f(z)=z+\sum_{n\geq2}a_nz^n$ on $\mathbb D$. Let $\mathcal K\subset\mathcal A$ consist of univalent functions with convex image. For analytic series $u=\sum u_nz^n$ and $v=\sum v_nz^n$, their Hadamard convolution is $(u*v)(z)=\sum u_nv_nz^n$. Ruscheweyh's class in the source is
\[
 M=\{f\in\mathcal A:f*\phi\text{ is univalent for every }\phi\in\mathcal K\}.
\]
The source reports that all close-to-convex functions belong to $M$.

Let $\mathcal P=\{F:F\text{ analytic on }\mathbb D,\ F(0)=1,\ \operatorname{Re}F>0\}$. This is the usual normalization for positive-real-part functions. Define the displayed condition on $g\in\mathcal A$ by
\[
 \operatorname{Re}\left\{\frac{\phi*(gF)}{\phi*g}\right\}>0
 \quad\text{on }\mathbb D,
 \qquad\forall\phi\in\mathcal K\ \forall F\in\mathcal P. \tag{1}
\]
The product $gF$ is ordinary multiplication before taking convolution. At $z=0$ the quotient has removable value $1$. At any other possible denominator zero the intended analytic interpretation must be specified; the source does not state that convention. Let $\mathcal G$ denote the functions satisfying (1) with that issue resolved.

\textbf{First question (T. Sheil-Small).} The source says (1) holds for starlike $g$, meaning normalized univalent functions whose images are starlike about $0$. Does $\mathcal G$ contain any nonstarlike functions, or a significant larger family than the starlike class? The existence target is $\mathcal G\setminus S^*\ne\varnothing$, with a larger-class description also requested.

\textbf{Classification question.} The source next states that if $g$ satisfies a condition referenced as ``6.2'' and
\[
 \operatorname{Re}\left(\frac{zf'(z)}{g(z)}\right)>0\quad(z\in\mathbb D), \tag{2}
\]
then $f\in M$, and asks whether this classifies $M$. If the missing reference ``6.2'' means condition (1), the natural converse question is
\[
 \forall f\in M\quad\exists g\in\mathcal G:
 \operatorname{Re}\left(\frac{zf'}g\right)>0\text{ on }\mathbb D.
\]
The quotient in (2) also has removable value $1$ at $0$. The record uses inconsistent source labels C6.1 and 6.2; identifying the missing reference with (1) is conditional, not an authenticated correction.

The supplied summary records OPEN-TRIAGE, says the item was open in the 2018 edition, and reports no definitive resolution located. The reported sufficient conditions remain attributed source statements.
\subsection{Short English statement}
Describe which normalized analytic functions satisfy the universal positive-real-part convolution condition beyond starlike functions. Then ask whether the associated inequality $\operatorname{Re}(zf'/g)>0$ characterizes Ruscheweyh's class $M$. A broken condition reference and quotient conventions need clarification.
\subsection{Sources}
\begin{itemize}
\item Local UnsolvedMath registry, permanent ID 2306034, source code catalogue entry, ``Research Problems in Function Theory --- Problem 6.34''; attribution: T. Sheil-Small. Supplied status open, classification OPEN-TRIAGE, updated 2026-08-06.
\item Hayman, \emph{Research Problems in Function Theory} (2018), Problem 6.34, \url{https://arxiv.org/abs/1809.07200}. The record reports source-TeX extraction and access on 2026-07-29; the 2018 update was retained for status triage. Current status and rights require review; the record describes publicly available arXiv source TeX.
\item Ruscheweyh's class and the reported sufficient convolution conditions are named in the source without a separate bibliography. The source labels C6.1 and 6.2 are retained as unresolved cross-references.
\end{itemize}
\end{document}
